% Part opener. Structural figures get their own number prefix so they
% do not inherit a misleading chapter number.
\structuralfigures{P6.}

\vspace*{0.4cm}
Everything up to here has been about networks. These last two chapters are about
the exam --- one on the method that the largest group of topics is really
testing, and one on sitting the paper itself.

They are short, and they are not filler. The single most common way to fail an
exam you know the material for is to mismanage it: to spend eleven minutes on a
question worth the same as one that takes forty seconds, or to reach a
configuration task with no time left, or to answer what you assumed was asked
rather than what was.

\ccnafigh{%
\begin{tikzpicture}[font=\scriptsize]
  \tikzset{
    seg/.style={draw=#1, fill=#1!16, text=#1!40!black, minimum height=1.15cm,
                align=center, font=\scriptsize\bfseries, line width=0.9pt},
    lb/.style={font=\tiny, text=neutral!85!black, align=center, anchor=north}}

  % 29 topics scaled to 14.4 cm: 0.4966 cm per topic.
  \node[seg=dom4, minimum width=7.95cm] (a) at (3.97,0)  {16 topics\\[1pt]\Large 55\%};
  \node[seg=dom5, minimum width=3.48cm] (b) at (9.69,0)  {7\\[1pt]\Large 24\%};
  \node[seg=dom1, minimum width=2.98cm] (c) at (12.92,0) {6\\[1pt]\Large 21\%};

  \node[lb, text width=7.8cm] at (3.97,-0.72)
      {\textbf{Configure, Troubleshoot, Diagnose}\\
       do it, or fix it, at a command line};
  \node[lb, text width=3.4cm] at (9.69,-0.72)
      {\textbf{Interpret, Validate,\\Select, Use, Manage}\\read output and judge it};
  \node[lb, text width=2.9cm] at (12.92,-0.72)
      {\textbf{Describe}\\recall and explain};
\end{tikzpicture}}%
{The twenty-nine topics by their opening verb. Only a fifth of the blueprint asks
you to recall anything; the rest asks you to do something or to read something
and decide.}{verbs}

\begin{examnote}[The number that should shape your last month of revision]
Count the verbs in Cisco's own topic list and the distribution above falls out:
eight topics begin \emph{Configure}, six \emph{Troubleshoot}, two
\emph{Diagnose}. That is sixteen of twenty-nine --- \textbf{55\%} --- where
knowing the answer is not enough, because you have to produce it at a CLI or find
it in someone else's broken configuration.

Revision that consists of re-reading notes prepares you for the 21\%. The other
79\% is prepared for by building things and breaking them, which is why every
chapter of this volume has a lab and a Break \& Fix rather than a summary and a
quiz.

If you have four weeks left and have read everything once, spend three of them at
a command line.
\end{examnote}
